%-------------------------------------------------------------------------------
% Mattia Giovannini — CV en français (version publique), mise en page Awesome-CV
% Compiler avec LuaLaTeX :  latexmk -lualatex CV_FR-public.tex
%-------------------------------------------------------------------------------
\documentclass[11pt, a4paper]{../awesome-cv}
% Shared Awesome-CV settings for the public CVs (IT/EN/FR).
% Each CV loads this after \documentclass{../awesome-cv}; build with LuaLaTeX.

\geometry{left=1.4cm, top=.9cm, right=1.4cm, bottom=1.6cm, footskip=.5cm}

% Same blue accent as the other CVs
\definecolor{awesome}{RGB}{80, 100, 200}
\setbool{acvSectionColorHighlight}{true}
\renewcommand{\acvHeaderSocialSep}{\quad\textbar\quad}

% No hyphenation: words stay whole when application portals extract the text
\hyphenpenalty=10000
\exhyphenpenalty=10000

\name{Mattia}{Giovannini}
\email{mattia.giovannini6@studio.unibo.it}
\homepage{mattia-giovannini.github.io}
\github{frigori}

% Re-load the class fonts with LuaTeX's node renderer: with HarfBuzz, explicit
% hyphens ("eye-tracking") are extracted from the PDF as invisible soft hyphens.
\defaultfontfeatures{Ligatures=TeX, Renderer=Node}
\setmainfont{Source Sans 3}[
  UprightFont=*, ItalicFont=* Italic, BoldFont=* Bold, BoldItalicFont=* Bold Italic,
  FontFace={l}{n}{Font=* Light}, FontFace={l}{it}{Font=* Light Italic},
]
\setsansfont{Source Sans 3}[
  UprightFont=*, ItalicFont=* Italic, BoldFont=* Bold, BoldItalicFont=* Bold Italic,
  FontFace={l}{n}{Font=* Light}, FontFace={l}{it}{Font=* Light Italic},
]
\renewfontfamily\roboto{Roboto}[
  UprightFont=*, ItalicFont=* Italic, BoldFont=* Bold, BoldItalicFont=* Bold Italic,
  FontFace={l}{n}{Font=* Light}, FontFace={l}{it}{Font=* Light Italic},
]


\position{Étudiant en Master d'ingénierie informatique (double diplôme){\enskip\cdotp\enskip}Licence en ingénierie informatique}
\address{Bologne, Italie}

\begin{document}

\makecvheader[C]
\makecvfooter{}{Mattia Giovannini~~~·~~~Curriculum Vitae}{\thepage}

%-------------------------------------------------------------------------------
\cvsection{Profil}

\begin{cvparagraph}
Diplômé en ingénierie informatique à l'Université de Bologne et inscrit en double diplôme : Master en ingénierie
informatique (Université de Bologne) et Master Ubiquitous Systems and Networks (Université Côte d'Azur, Polytech Nice
Sophia), avec un intérêt marqué pour l'intelligence artificielle, le machine learning et l'analyse de données. Ma
formation m'a donné des bases solides en algorithmique, statistiques, systèmes et génie logiciel, que j'ai appliquées
dans un contexte de recherche lors d'un semestre Erasmus en France. Je suis motivé par la conception de systèmes
intelligents pour résoudre des problèmes concrets.
\end{cvparagraph}

%-------------------------------------------------------------------------------
\cvsection{Formation}
\begin{cventries}
  \cventry
    {Double diplôme avec le Master Ubiquitous Systems and Networks, Université Côte d'Azur}
    {Master en ingénierie informatique — double diplôme}
    {Bologne, Italie \& Nice, France}
    {2026 -- 2028 (en cours)}
    {}

  \cventry
    {Note finale : 94/110}
    {Licence en ingénierie informatique — Université de Bologne}
    {Bologne, Italie}
    {2021 -- juillet 2026}
    {
      \begin{cvitems}
        \item {\textbf{Mémoire :} \textit{Détection du trouble du spectre de l'autisme à partir de données d'eye-tracking par machine learning}, classification et analyse de signaux appliquées à des enregistrements cliniques d'oculométrie.}
        \item {Cours pertinents : analyse mathématique, algèbre linéaire et géométrie, probabilités et statistiques, algorithmique, systèmes d'exploitation, réseaux informatiques, systèmes d'information, génie logiciel, technologies web, architecture des ordinateurs.}
      \end{cvitems}
    }

  \cventry
    {Projet d'études Erasmus}
    {Université de Tours}
    {Tours, France}
    {2025}
    {
      \begin{cvitems}
        \item {Projet sur le \textbf{machine learning appliqué à l'analyse du trouble du spectre de l'autisme (TSA)} : classification et extraction de caractéristiques sur des données d'eye-tracking.}
        \item {Travail dans un environnement académique international et francophone ; renforcement en analyse de données, Python et méthodologie scientifique.}
      \end{cvitems}
    }

  \cventry
    {Diplôme de lycée scientifique}
    {Liceo Augusto Righi}
    {Bologne, Italie}
    {2016 -- 2021}
    {}
\end{cventries}

%-------------------------------------------------------------------------------
\cvsection{Projets}
\begin{cventries}
  \cventry
    {Équipe projet GenAI}
    {BOOM GEN AI Challenge 2026 — VEM Sistemi}
    {}
    {2026}
    {
      \begin{cvitems}
        \item {Contribution au développement d'un \textbf{assistant GenAI} pour l'analyse automatisée d'appels d'offres et l'aide à la décision \textbf{GO/NO-GO}.}
        \item {Conception et implémentation d'un \textbf{module de scoring en Python} pour évaluer l'alignement entre les exigences d'un appel d'offres et les capacités de l'entreprise.}
        \item {\textbf{Prompt engineering} et services \textbf{Azure OpenAI} pour l'extraction, la classification et l'enrichissement des exigences ; analyse d'un pipeline de parsing documentaire.}
      \end{cvitems}
    }

  \cventry
    {Project work · Enterprise Platforms and DevOps M, Université de Bologne}
    {Scheduler Kubernetes sensible aux ressources}
    {}
    {2026 (en cours)}
    {
      \begin{cvitems}
        \item {Développement d'un \textbf{scheduler Kubernetes personnalisé} en Python qui place chaque pod sur le nœud de plus faible latence disposant de CPU, mémoire et capacité suffisantes.}
        \item {Mesure de la latence des nœuds en temps réel, binding via l'API Kubernetes et justification de chaque décision par des événements \textit{Scheduled}/\textit{FailedScheduling}.}
        \item {Extension en cours vers le \textbf{placement intelligent et le remplacement} de services sur des nœuds edge mobiles (robots humanoïdes entre postes de travail).}
      \end{cvitems}
    }

  \cventry
    {Discrete and Network Optimization Problems M, Université de Bologne}
    {Benchmark de restauration après panne de lien SDN}
    {\href{https://github.com/frigori/net-restoration-bench}{\faGithub\ net-restoration-bench}}
    {2026 (en cours)}
    {
      \begin{cvitems}
        \item {Comparaison de la \textbf{protection par chemins de secours partagés (SBPP)}, du \textbf{Fast Reroute} et du \textbf{recalcul global MILP} : surcoût en capacité, taux de restauration, temps de réaction.}
        \item {Python avec NetworkX et PuLP ; une interface plan/react commune exécute les stratégies sur les mêmes topologies et scénarios de panne.}
      \end{cvitems}
    }

  \cventry
    {Lean Startup Lab · Technology Entrepreneurship}
    {AgenticOS — Audit IA multi-agents}
    {\href{https://agenticos-mvp.netlify.app}{\faGlobe\ MVP en ligne}}
    {2026}
    {
      \begin{cvitems}
        \item {Conception d'un \textbf{système multi-agents Red Team / Blue Team} (Gemini) qui explique pourquoi les agents IA d'achat ignorent un produit de mode.}
        \item {Un agent Blue Team compile les correctifs en \textbf{JSON-LD} et \textbf{llms.txt} et les évalue en sandbox ; MVP avec interface web et moteur conteneurisé.}
      \end{cvitems}
    }

  \cventry
    {Projet personnel}
    {Homelab auto-hébergé}
    {\href{https://mattia-giovannini.github.io/homelab/}{\faGlobe\ page homelab}}
    {En cours}
    {
      \begin{cvitems}
        \item {Exploitation de services conteneurisés sur des nœuds Raspberry Pi avec \textbf{Docker, Podman rootless et Kubernetes}, derrière un reverse proxy avec TLS.}
        \item {Supervision Prometheus/Grafana, sécurité CrowdSec ; banc d'essai pour donner à l'IA agentique un \textbf{accès limité et auditable}.}
      \end{cvitems}
    }
\end{cventries}

%-------------------------------------------------------------------------------
\cvsection{Compétences}
\begin{cvskills}
  \cvskill{Langages}{C, Java, Python, JavaScript, HTML/CSS, SQL, Bash}
  \cvskill{ML \& données}{scikit-learn, NumPy, pandas, Matplotlib, Jupyter Notebook}
  \cvskill{Systèmes \& réseaux}{Linux, TCP/IP, programmation socket, bases de données relationnelles}
  \cvskill{Outils}{Git, \LaTeX, MATLAB, Docker, Podman, Kubernetes}
  \cvskill{Centres d'intérêt}{Machine learning, analyse de données, IA générative, conteneurisation, automatisation système}
\end{cvskills}

%-------------------------------------------------------------------------------
\cvsection{Langues}
\begin{cvskills}
  \cvskill{Italien}{Langue maternelle}
  \cvskill{Anglais — C1}{Duolingo English Test 130/160 (C1) · CLA C1, Université de Bologne · Cambridge First Certificate B2 (178/190)}
  \cvskill{Français — C1}{CLA C1, Université de Bologne · cours B2/C1 suivis à l'Université de Tours}
\end{cvskills}

%-------------------------------------------------------------------------------
\cvsection{Autres expériences}
\begin{cventries}
  \cventry
    {Webmaster et rédacteur étudiant}
    {Liceo Augusto Righi}
    {Bologne, Italie}
    {2020 -- 2021}
    {
      \begin{cvitems}
        \item {Conception et gestion de plateformes étudiantes (\textit{Righi Network} et \textit{Il Sottobanco}), au service de la communication numérique et du journal de l'école.}
      \end{cvitems}
    }

  \cventry
    {Animateur bénévole en centre d'été}
    {Estate Ragazzi}
    {Bologne, Italie}
    {2016 -- 2020}
    {
      \begin{cvitems}
        \item {Gestion de groupe, travail en équipe, communication et responsabilité dans des activités éducatives.}
      \end{cvitems}
    }
\end{cventries}

\end{document}
